% !TEX root = ../main.tex
\chapter{Control espacial de qualitat}\label{ch:control}
\fancyhead[RE]{\slshape \textsf{Capítol \thechapter.\ Control espacial de qualitat}}
\fancyhead[LO]{\slshape \textsf{\nouppercase \rightmark}}

\section{De lambda global a Q-map}\label{sec:qmap_intro}

Els capítols anteriors han definit la relació entre qualitat, modulació i
latents. Aquest capítol fixa com es converteix aquesta idea en una decisió de
disseny: en lloc d'utilitzar una única qualitat per tota la imatge, el sistema
pot assignar una qualitat diferent a cada cel·la latent. En la ruta C, aquesta
informació és un Q-map de $32\times32$ bytes. Cada posició correspon a una
cel·la latent associada a una regió de control de $16\times16$ píxels en una
imatge $512\times512$.

El control espacial es construeix amb tres capes conceptuals:

\begin{enumerate}[leftmargin=*]
  \item una \textbf{base de qualitat}, que defineix el valor inicial de Q abans
  de considerar cap regió d'interès;
  \item un \textbf{preset}, que calcula una màscara binària sobre els blocs i
  indica quins es tractaran com a ROI;
  \item una \textbf{política}, que decideix com transformar la base i la màscara
  en el Q final de cada bloc.
\end{enumerate}

La base respon a la pregunta \emph{de quin nivell de qualitat partim?}. Pot ser
uniforme, si tots els blocs comencen amb el mateix Q, o adaptativa, si cada bloc
rep un valor inicial diferent segons el seu comportament local. El preset respon
a la pregunta \emph{quins blocs són importants per aquest criteri?}. La política
respon a la pregunta \emph{què fem amb els blocs marcats i amb els no marcats?}

Un exemple abstracte és:

\begin{equation}
\text{base de qualitat}
\;+\;
\text{màscara ROI}
\;+\;
\text{regla d'assignació}
\;\longrightarrow\;
\text{Q-map final}.
\end{equation}

El compressor no interpreta si una regió és vegetació, aigua, núvol o contrast
visible. Només rep el Q-map final. Aquesta separació és important perquè permet
canviar el criteri que detecta la ROI sense modificar el còdec, i permet canviar
la política de qualitat sense reescriure el detector.

\section{Escala de lambda i calibració}\label{sec:lambda005}

\subsection{Escala inicial del treball previ}

El treball anterior utilitzava una escala global amb $\lambda_{\max}=0,125$.
En aquella configuració, el punt històric $\lambda=0,1$ es quantitzava com:

\begin{equation}
Q=\operatorname{round}\left(255\frac{0,1}{0,125}\right)=204.
\end{equation}

Aquesta relació era suficient per reproduir una qualitat global fixada a tota
la imatge. En passar a control espacial, però, la decisió ja no és un únic valor:
cal indicar un Q diferent per cel·la latent, reconstruir-lo al descompressor i
comprovar que la qualitat demanada és la qualitat aplicada. Per això convé fixar
una escala explícita i única abans de comparar qualitat local i mida codificada.

\subsection{Escala utilitzada en aquesta memòria}

Per aquest motiu es defineix una escala explícita:

\begin{equation}
\lambda(Q)=\frac{Q}{255}\cdot0,05.
\label{eq:lambda005_control}
\end{equation}

Aquesta configuració s'anomena \code{lambda005}. El nom només indica el màxim
de l'escala de lambda. No canvia l'arquitectura del model ni els pesos
entrenats; canvia la manera com es tradueix un enter Q de 8 bits al valor de
$\lambda$ que rep el modulador.

La conseqüència pràctica és que qualsevol Q-map queda definit per una cadena
única:

\begin{equation}
Q_i \longrightarrow \lambda_i \longrightarrow M(\lambda_i)
\longrightarrow \text{quantització dels latents}.
\end{equation}

Els capítols de metodologia i resultats comproven després que aquesta cadena es
manté coherent en els experiments. Aquí només es fixa la convenció que utilitza
la resta de la memòria.

\subsection{Calcular, decidir i validar}

La idea principal del control fixed-quality és calcular quina $\lambda$, i per
tant quina Q, necessita una imatge o una regió per acostar-se a una qualitat
objectiu. Per fer-ho no n'hi ha prou amb mirar el preset: cal conèixer l'error
base que produeix el model sobre aquella imatge. En la ruta de càlcul que
s'avalua en aquest treball, aquest error base s'obté amb el còdec real:

\begin{enumerate}[leftmargin=*]
  \item \textbf{Compressió i descompressió base}. La imatge es comprimeix amb
  una qualitat de referència i es descomprimeix per obtenir una reconstrucció.
  Aquesta passada permet mesurar MSE o PSNR global i per bloc.
  \item \textbf{Càlcul de $\lambda$/Q}. L'error mesurat s'introdueix a les
  fórmules fixed-quality. El resultat és una $\lambda$ objectiu, que es quantitza
  com a Q i s'escriu al Q-map.
  \item \textbf{Aplicació del Q-map}. El Q-map calculat es combina amb la base,
  el preset i la política, i s'aplica en una compressió final. Aquesta etapa dona
  la reconstrucció final, el PSNR, els temps i, amb CSMR, el bitrate codificat.
\end{enumerate}

\begin{figure}[H]
\centering
\fbox{\begin{minipage}{0.92\textwidth}
\centering
Imatge
$\rightarrow$ compressió/descompressió base
$\rightarrow$ MSE/MSE$_0$
$\rightarrow$ $\lambda$/Q
$\rightarrow$ Q-map
$\rightarrow$ compressió final
$\rightarrow$ mètriques reals
\end{minipage}}
\caption{Ruta principal de càlcul de qualitat: el còdec real proporciona l'error
base, les fórmules fixed-quality dedueixen $\lambda$/Q i el Q-map resultant
s'aplica a la compressió final.}
\label{fig:calibration_decision_validation}
\end{figure}

Aquest plantejament explica per què el compressor i el descompressor C són tan
importants. No són només eines de comprovació posterior: formen part del cost que
cal avaluar si es vol calcular qualitat a bord amb error mesurat. Una alternativa
pràctica és guardar relacions ja mesurades en una taula auxiliar. Aquesta idea
serveix per accelerar campanyes, comparar polítiques i mantenir traçabilitat. Els
capítols següents separen aquests dos usos: el càlcul amb error mesurat i la
taula com a suport experimental.

\begin{table}[H]
\centering
\scriptsize
\begin{tabularx}{\textwidth}{p{2.6cm}X X X}
\toprule
Estratègia & Com decideix Q & Cost a bord & Ús en aquesta memòria \\
\midrule
Càlcul amb error mesurat & Comprimeix i descomprimeix una qualitat base, mesura
MSE/MSE$_0$ i aplica les fórmules fixed-quality per deduir $\lambda$/Q. & Alt:
inclou una passada de compressor i descompressor abans o dins la decisió de
qualitat. & Ruta principal que es vol avaluar per processament a bord. Justifica
les mesures de temps i les optimitzacions C. \\
Taula precalibrada & Reutilitza errors i coeficients mesurats prèviament per
calcular Q sense repetir la passada base en cada comparativa. & Baix: lectura de
taula, càlcul del preset i escriptura del Q-map. & Suport pràctic per accelerar
experiments, comparar polítiques, fixar rangs i reproduir resultats. \\
Q fixa directa & Assigna valors decidits sense model predictiu, per exemple
Q uniforme o ROI/fons fixos. & Molt baix. & Base de \code{q204} i
\code{preserve-roi}; actua com a control i com a política de conservació de ROI. \\
\bottomrule
\end{tabularx}
\caption{Estratègies per decidir la qualitat. La ruta amb error mesurat és la
lectura principal del treball; la taula precalibrada és una eina auxiliar per fer
comparatives reproduïbles sense repetir sempre la passada base.}
\label{tab:q_decision_strategies}
\end{table}

\subsection{Calibració per bloc}

La calibració per bloc és el pont entre les fórmules del capítol anterior i les
decisions discretes que ha de prendre el sistema. En la ruta amb error mesurat,
el MSE base prové de la passada real del còdec. En les campanyes àmplies, aquest
mateix coneixement es pot guardar en una taula per evitar repetir la mesura en
cada comparativa. Per això la taula no és la contribució central, sinó una
representació compacta de mesures i ajustos previs.

La taula auxiliar apareix com una resposta pràctica al cost de repetir sempre la
compressió i la descompressió base. La idea és mesurar prèviament diversos valors
discrets de Q, observar l'error per bloc i ajustar el model:

\begin{equation}
\mathrm{MSE}\simeq c_0+\frac{c_1}{(m_a\lambda+m_b)^2}.
\end{equation}

Per cada posició de la graella latent, aquesta representació guarda coeficients
ajustats, rangs de Q disponibles i errors observats en punts de referència. Amb
aquesta informació, els generadors de Q-map poden aproximar decisions de qualitat
de forma ràpida. Els capítols de metodologia i resultats expliquen després com
s'ha utilitzat aquesta eina i fins a quin punt els Q-maps generats mantenen la
qualitat prevista quan s'apliquen al còdec real.

Com a eina auxiliar, la taula permet als generadors de Q-map fer tres operacions:

\begin{itemize}[leftmargin=*]
  \item convertir un objectiu global de qualitat en una Q constant;
  \item construir una base variable mantenint una mitjana controlada;
  \item assignar valors fixos de Q a ROI i fons quan la política ho requereix.
\end{itemize}

Si es canvien pesos, nombre de bandes, normalització o escala de lambda, aquesta
taula auxiliar ha de revisar-se. Altrament, el Q-map continuaria sent
sintàcticament vàlid, però podria deixar de representar la qualitat prevista. En
canvi, la ruta amb error mesurat torna a obtenir el MSE base amb el còdec i és la
que fixa la interpretació física del càlcul de $\lambda$/Q.

\section{Rang de Q}\label{sec:qrange}

El Q-map es desa com un enter sense signe de 8 bits. Per construcció, això
permet:

\begin{equation}
0\leq Q\leq255.
\end{equation}

El valor alt de Q representa més qualitat i el valor baix representa més
degradació. La configuració principal d'aquesta memòria utilitza:

\begin{equation}
128\leq Q\leq255.
\end{equation}

Aquest interval no és una propietat física universal del model. És el rang que
s'utilitza com a base de treball per mantenir el control de qualitat dins la
zona més estudiada de la calibració. Per sota de Q128 el format continua sent
representable i pot ser útil quan es vol degradar més agressivament el fons.
Tanmateix, aquesta extensió requereix una anàlisi pròpia perquè poden aparèixer
tres efectes:

\begin{itemize}[leftmargin=*]
  \item la relació entre Q i error pot dependre més del contingut concret del
  bloc;
  \item un objectiu de PSNR pot saturar al límit de la taula i quedar més lluny
  del valor demanat;
  \item la degradació pot afectar visualment zones properes per l'efecte de
  veïnatge de les convolucions.
\end{itemize}

Per tant, la memòria separa la configuració principal, basada en Q128--Q255, de
les proves amb Q més baixes que s'analitzen més endavant com a línia de recerca.

\section{Modes de Q-map}\label{sec:modes}

Un mode de Q-map és una regla que omple les 1.024 posicions de la graella
latent. La Taula~\ref{tab:modes} resumeix els modes que estructuren la memòria.
Els modes auxiliars de desenvolupament queden fora del cos principal perquè no
són necessaris per entendre la comparació final.

\begin{table}[H]
\centering
\scriptsize
\begin{tabularx}{\textwidth}{p{2.5cm}p{3.2cm}X p{2.5cm}}
\toprule
Mode & Tipus de decisió & Q resultant & Paper dins la memòria \\
\midrule
\code{q204} & Qualitat uniforme explícita & Tots els blocs reben Q=204. &
Referència uniforme principal. \\
\code{global\_target} & Objectiu global de MSE o PSNR & Tots els blocs reben el
mateix Q deduït de la calibració. & Connexió directa amb qualitat fixa. \\
\code{adaptive\_s8} & Dificultat relativa per bloc & Q variable amb mitjana
propera a 204. & Control tècnic sense ROI. \\
\code{focus\_bgq128} & Màscara ROI i fons degradat & ROI amb base adaptativa
reforçada; fons a Q128. & Política principal de redistribució espacial. \\
\code{preserve-roi} & Màscara ROI i Q explícita & ROI a Q alta fixa; fons a
Q128. & Política orientada a conservar més la ROI. \\
\bottomrule
\end{tabularx}
\caption{Modes de Q-map utilitzats a la memòria. La taula descriu la regla de
construcció, no els resultats obtinguts.}
\label{tab:modes}
\end{table}

\subsection{q204 i global-target}

\code{q204} assigna el mateix valor a tots els blocs:

\begin{equation}
Q_i=204\quad\forall i.
\end{equation}

És una ordre simple: es demana explícitament Q=204 i no es calcula cap objectiu
d'error. Per això és una referència fàcil d'interpretar.

\code{global\_target} també genera un mapa constant, però l'ordre d'usuari és
diferent. En lloc de donar Q directament, es dona un objectiu de MSE o PSNR i la
calibració en dedueix el valor:

\begin{equation}
Q_i=Q^\star(\mathrm{target})\quad\forall i.
\end{equation}

Els dos modes poden donar el mateix Q si l'objectiu condueix a 204, però no
responen a la mateixa pregunta. \code{q204} pregunta \emph{què passa amb aquesta
Q fixa?}; \code{global\_target} pregunta \emph{quina Q necessito per aquest
objectiu global?}

\subsection{adaptive-s8}

\code{adaptive\_s8} no detecta classes físiques ni regions semàntiques. Només
redistribueix Q segons la dificultat relativa del bloc. En la ruta amb error
mesurat, aquesta dificultat surt del MSE base obtingut amb la passada de
compressió i descompressió de referència. En les campanyes comparatives de la
memòria, el mateix valor es reutilitza des de la taula auxiliar de calibració per
evitar repetir la mesura per cada variant.

Per convertir aquest error en una dificultat relativa es treballa en escala
logarítmica:

\begin{equation}
z_i=
\frac{\log(\mathrm{MSE}_i)-\mu_{\log\mathrm{MSE}}}
{\sigma_{\log\mathrm{MSE}}}.
\end{equation}

Aquí $\mu_{\log\mathrm{MSE}}$ i $\sigma_{\log\mathrm{MSE}}$ normalitzen els errors
de la mateixa imatge o del mateix cas calibrat. Per tant, $z_i$ no diu si un bloc
és difícil de manera universal ni analitza semàntica; diu si aquell bloc té més
o menys error que la mitjana del cas. Això és útil perquè manté una Q mitjana
controlada mentre redistribueix qualitat. La limitació és que \code{adaptive\_s8}
per si sol no sap on hi ha vegetació, núvols o aigua: és una base de qualitat, no
un preset.

La regla utilitzada és:

\begin{equation}
Q_i=\operatorname{clip}(204+8z_i).
\label{eq:adaptive_s8}
\end{equation}

La funció $\operatorname{clip}$ limita el valor calculat perquè quedi dins del
rang de Q permès. El centre 204 prové de la referència històrica del projecte i
el factor 8 controla la intensitat de la variació. Amb un factor petit, el mapa
s'assembla molt a una Q constant; amb un factor gran, molts blocs poden saturar
als límits.

Dit d'una altra manera, \code{adaptive\_s8} és viable a bord perquè és una regla
de consulta i normalització sobre una taula ja disponible. No és un detector de
ROI i no substitueix els presets automàtics que es defineixen més avall.

\subsection{focus-bgq128}

\code{focus\_bgq128} és la primera política que combina les dues peces: una base
adaptativa i una màscara ROI. La màscara $m_i$ la proporciona un preset; $m_i=1$
indica ROI i $m_i=0$ indica fons. La base $Q_{\mathrm{adaptive},i}$ és el valor
que hauria produït \code{adaptive\_s8} per aquell bloc. La regla és:

\begin{equation}
Q_i=
\begin{cases}
  \operatorname{clip}(Q_{\mathrm{adaptive},i}+16),&m_i=1,\\
128,&m_i=0.
\end{cases}
\label{eq:focus}
\end{equation}

La ROI no rep una Q fixa. Cada bloc ROI conserva la variació de la base
adaptativa i rep un increment de 16 unitats de Q. Això vol dir que dos blocs dins
la mateixa ROI poden acabar amb Q diferents si la seva dificultat calibrada era
diferent. El fons, en canvi, deixa de seguir la base adaptativa i queda fixat a
Q128.

Així, \code{adaptive\_s8} i els presets no són alternatives incompatibles.
\code{adaptive\_s8} tot sol és un mode sense ROI; \code{preset + focus\_bgq128}
és una combinació on el preset decideix quins blocs són ROI i la política
\code{focus\_bgq128} utilitza \code{adaptive\_s8} només com a base per a aquests
blocs ROI. El compressor final no veu aquesta descomposició: només rep el Q-map
ja calculat.

\subsection{preserve-roi}

\code{preserve-roi} utilitza la mateixa idea de màscara, però canvia la regla
d'assignació dins la ROI:

\begin{equation}
Q_i=
\begin{cases}
Q_{\mathrm{ROI}},&m_i=1,\\
Q_{\mathrm{fons}},&m_i=0,
\end{cases}
\quad
Q_{\mathrm{ROI}}\in\{240,255\},\quad Q_{\mathrm{fons}}=128.
\label{eq:preserve}
\end{equation}

La diferència essencial és que \code{focus\_bgq128} manté una ROI adaptativa,
mentre que \code{preserve-roi} imposa el mateix Q alt a tots els blocs ROI. Si
la prioritat és protegir qualsevol bloc marcat com a ROI, \code{preserve-roi}
és més conservador. Si la prioritat és evitar donar qualitat alta a blocs ROI
que no la necessiten tant, \code{focus\_bgq128} és més selectiu.

\section{Presets automàtics}\label{sec:presets}

Un preset és una regla lleugera que transforma la imatge en una màscara ROI. No
és una xarxa neuronal ni una classificació supervisada completa. Calcula una
magnitud per píxel o per bloc, la compara amb un llindar i marca els blocs que
compleixen el criteri. La Taula~\ref{tab:index_presets} presenta tots els
presets considerats abans d'entrar en les famílies de càlcul.

\begin{table}[H]
\centering
\small
\begin{tabularx}{\textwidth}{p{3.0cm}p{2.5cm}p{2.0cm}X}
\toprule
Preset & Magnitud & Llindar & Lectura de la màscara ROI \\
\midrule
\code{vegetation} & NDVI B08/B04 & $\geq0,40$ & Blocs amb resposta de
vegetació moderada o alta. \\
\code{high\_ndvi} & NDVI B08/B04 & $\geq0,50$ & Blocs amb resposta NDVI més
restrictiva. \\
\code{low\_ndvi} & NDVI B08/B04 & $\leq0,15$ & Blocs amb NDVI baix; pot incloure
sòl, cobertes no vegetades o valors negatius. \\
\code{vegetation\_green} & GNDVI B08/B03 & $\geq0,50$ & Blocs amb vegetació
contrastada amb la banda verda. \\
\code{chlorophyll} & NDCI B05/B04 & $\geq0,10$ & Blocs amb resposta associada a
clorofil·la segons les bandes disponibles. \\
\code{water\_body} & NDWI B03/B08 & $\geq0,10$ & Blocs compatibles amb presència
d'aigua segons l'índex utilitzat. \\
\code{dark\_regions} & Mitjana visible & $\leq0,26$ & Blocs foscos en les bandes
visibles. \\
\code{local\_contrast} & Desviació visible & $\geq0,035$ & Blocs amb variació
visible alta. \\
\code{clouds} & Fracció CBY & $\geq0,50$ & Blocs amb alta fracció de píxels
candidats a núvol. \\
\code{cloud\_avoid} & Fracció CBY & $\leq0,50$ & Blocs que eviten la màscara de
núvol segons el detector disponible. \\
\bottomrule
\end{tabularx}
\caption{Presets automàtics considerats. Els llindars són paràmetres de treball
per generar màscares ROI; no són constants universals de teledetecció.}
\label{tab:index_presets}
\end{table}

\subsection{Índexs espectrals normalitzats}

Els presets basats en dues bandes utilitzen una forma normalitzada:

\begin{equation}
I(A,B)=\frac{A-B}{A+B}.
\label{eq:index}
\end{equation}

El valor es calcula per píxel vàlid i després es resumeix dins cada bloc. NDVI i
NDWI tenen fonament bibliogràfic \cite{rouse1974ndvi,mcfeeters1996ndwi}, però
els valors de llindar depenen de l'escena, la resposta radiomètrica i les bandes
disponibles. Per això, en aquesta memòria el llindar s'ha de llegir com una
decisió de configuració, no com una definició física tancada.

Aquesta precaució és especialment important en \code{low\_ndvi}. En NDVI,
valors positius elevats solen associar-se a vegetació activa; valors propers a
zero poden correspondre a sòl nu o cobertes urbanes; i valors negatius poden
aparèixer en aigua, ombres o superfícies amb baixa resposta NIR. Per tant,
\code{low\_ndvi} no s'interpreta com un detector urbà pur, sinó com un preset
ampli de resposta NDVI baixa que s'ha d'avaluar amb les màscares i resultats del
dataset.

\subsection{Contrast i brillantor visible}

Els presets visibles treballen amb les bandes B02, B03 i B04. Es calcula una
mitjana visible normalitzada:

\begin{equation}
V=\frac{B02+B03+B04}{3\cdot10000}.
\end{equation}

\code{dark\_regions} marca blocs amb valor visible mitjà baix.
\code{local\_contrast} marca blocs amb desviació estàndard visible alta. Aquests
criteris no identifiquen una classe física única; només defineixen una regió
d'interès segons brillantor o textura.

\subsection{Núvol i no-núvol amb les bandes disponibles}

El detector CBY bàsic utilitza bandes visibles:

\begin{align}
bRatio&=\frac{B03-0,175}{0,39-0,175},\\
NDGR&=\frac{B03-B04}{B03+B04}.
\end{align}

Un píxel és candidat a núvol si $bRatio>1$ o si $bRatio>0$ i $NDGR>0$. El valor
del bloc és la fracció de píxels candidats. Si B11/SWIR fos disponible es podria
afegir una condició espectral més informativa; amb el conjunt de vuit bandes
utilitzat en aquesta memòria s'aplica la versió visible. Això limita la lectura
física del preset, però no impedeix utilitzar-lo com a detector automàtic de
regions per construir Q-maps.

\section{Llindars i cobertura ROI}\label{sec:llindars}

Un llindar converteix una magnitud contínua en una màscara binària. Per cada
bloc $i$ es calcula un valor $s_i$ i s'aplica una regla del tipus:

\begin{equation}
m_i =
\begin{cases}
1,&s_i \geq \tau,\\
0,&s_i < \tau.
\end{cases}
\end{equation}

En alguns presets la desigualtat s'inverteix, com en el cas de NDVI baix o
regions fosques. El punt important és que el llindar no controla directament la
qualitat: només controla quins blocs entren a la ROI. La qualitat final depèn de
la política aplicada després.

La cobertura de ROI és el percentatge de blocs amb $m_i=1$:

\begin{equation}
C = 100\frac{\sum_i m_i}{1024}.
\end{equation}

Per interpretar aquesta cobertura s'utilitza la nomenclatura següent:

\begin{equation}
\begin{array}{ll}
\code{no\_roi}:&C=0\%,\\
\code{low\_roi}:&0<C<10\%,\\
\code{mid\_roi}:&10\%\leq C\leq70\%,\\
\code{high\_roi}:&70\%<C<100\%,\\
\code{all\_roi}:&C=100\%.
\end{array}
\label{eq:roi_groups_control}
\end{equation}

Aquestes etiquetes no són presets i no modifiquen cap Q-map. Són categories
d'anàlisi que ajuden a detectar casos límit. Si la ROI és buida, no hi ha regió
protegida. Si és total, no queda fons amb el qual redistribuir qualitat. Si és
molt petita o molt dispersa, el camp receptiu de la xarxa pot barrejar informació
amb blocs veïns i fer que la protecció sigui menys efectiva.

\section{Combinació preset + política}\label{sec:control_decision}

El sistema complet es pot llegir com una composició:

\begin{equation}
\text{preset}
\;+\;
\text{política}
\;\longrightarrow\;
\text{Q-map}.
\end{equation}

El preset només decideix quins blocs són ROI. La política només decideix quina Q
rep la ROI i quina Q rep el fons. El Q-map és l'únic objecte que consumeixen
compressor i descompressor. Aquesta separació permet comparar combinacions sense
canviar el model neuronal.

Per exemple, si un preset marca una zona central com a ROI, \code{focus\_bgq128}
donarà a aquesta zona una qualitat adaptativa reforçada i fixarà la resta a
Q128. Amb la mateixa màscara, \code{preserve-roi} donaria una Q alta uniforme a
tota la zona central i també fixaria el fons a Q128. El detector no canvia; el
que canvia és la política de qualitat.

\subsection{Traçabilitat de la implementació C}\label{sec:interface_c}

La implementació C reflecteix aquesta separació. Les eines de generació de Q-map
produeixen un fitxer binari de 1.024 bytes i un resum per bloc amb el valor del
preset, la decisió ROI/fons, la Q base i la Q final. El cos principal de la
memòria només necessita aquesta idea. Els noms exactes de paràmetres, ordres de
reproducció i fitxers de sortida es recullen als apèndixs i al repositori per no
interrompre el fil conceptual del capítol.

\subsection{Modes auxiliars}\label{sec:auxiliary_modes}

Durant el desenvolupament també es van implementar polítiques auxiliars, com
incrementar Q dins la ROI sense fixar el fons, o calcular un objectiu explícit
de qualitat només per la ROI. Aquestes variants són útils per depurar i entendre
el sistema, però la memòria es centra en les combinacions que responen a les
preguntes principals: qualitat uniforme, redistribució adaptativa, ROI protegida
amb fons degradat i ROI preservada amb Q alta.
